\section{Exponentiated Online-to-non-convex}
\label{sec:alg}

\begin{algorithm}[t]
\caption{Exponentiated O2NC}
\label{alg:O2NC-exp-avg}
\begin{algorithmic}[1]
    \STATE \textbf{Input:} OCO algorithm $\calA$, initial state $\vecx_0$, parameters $N\in\NN, \beta\in (0,1)$, regularizers $\calR_n(\Delta)$.
    \FOR {$n\gets 1,2,\ldots,N$}
        \STATE Receive $\Delta_n$ from $\calA$.
        \STATE Update $\vecx_n \gets \vecx_{n-1} + s_n\Delta_n$, where $s_n\sim\Expo(1)$ i.i.d.
        \STATE Compute $\vecg_n \gets \nabla f(\vecx_n, z_n)$.
        \STATE Send loss $\ell_n(\Delta) = \langle \beta^{-n}\vecg_n, \Delta\rangle + \calR_n(\Delta)$ to $\calA$.
        
        % \STATE 
        // For output only (does \textit{not} affect training):
        \STATE Update $\overline\vecx_n = \frac{\beta-\beta^n}{1-\beta^n}\overline\vecx_{n-1} + \frac{1-\beta}{1-\beta^n}\vecx_n$. \\
        Equivalently, $\overline \vecx_n = \sum_{t=1}^n \beta^{n-t}\vecx_t\cdot \frac{1-\beta}{1-\beta^n}$.
    \ENDFOR
    % \State Compute $\vecx^k \gets \frac{1}{T} \sum_{n\in I_k} \vecx_n$ where $I_k=[(k-1)T+1,kT]$.
    \STATE Output $\overline \vecx \sim \Unif(\{\overline\vecx_n:n\in[N]\})$.
\end{algorithmic}
\end{algorithm}



% \subsection{Algorithm Design}
% \label{sec:alg-design}

% In this section, we introduce a comprehensive framework designed to transform any OCO algorithm with sublinear regret into a stochastic optimization algorithm suitable for non-smooth non-convex losses. This framework is an adaptation of the Online-to-non-convex Conversion (O2NC) algorithm \cite{cutkosky2023optimal}. We have implemented several critical modifications to more effectively recover momentum algorithms. The specifics of each modification are elaborated in Section \ref{sec:alg-design}.

In this section, we present our improved online-to-non-convex framework, \alg, and explain the key techniques we employed to improve the algorithm. The pseudocode is presented in Algorithm \ref{alg:O2NC-exp-avg}.


\paragraph{Random Scaling}
% One notable difference between \alg \ and standard O2NC lies in the selection of the scaling random variable $s_n$. In O2NC, $s_n$ is chosen to be uniform, $s_n\sim \Unif([0,1])$, based on
% \begin{align*}
%     \E[F(\vecx+\Delta) - F(\vecx)] = \Ex_{s\sim\Unif([0,1])}[\langle \vecx+s\Delta), \Delta\rangle].
% \end{align*}
% Given $\vecx_n=\vecx_{n-1}+\Delta_n$ and $\vecw_n=\vecx_{n-1}+s_n\Delta_n$, it follows that $\langle \nabla F(\vecw_n),\Delta_n\rangle$ is an unbiased estimator of $F(\vecx_n)-F(\vecx_{n-1})$. However, this approach introduces additional memory burden, as it requires  to compute the gradient at the intermediate variable $\vecw_n$ lying between $\vecx_{n-1}$ and $\vecx_n$, rather than directly at the iterate $\vecx_n$.
% To address this limitation, we choose an exponential random variable as the scaling factor, namely $s_n \sim \Expo(1)$. The reason for this choice is as follows.

One notable feature of Algorithm \ref{alg:O2NC-exp-avg} is that the update $\Delta_n$ is scaled by an exponential random variable $s_n$. 
Formally, we have the following result:

\begin{restatable}{lemma}{ExponentialScaling}
\label{lem:exp-scaling}
Let $s\sim\mathrm{Exp}(\lambda)$ for some $\lambda>0$, then
\begin{equation*}
    \Ex_s[F(\vecx+s\Delta) - F(\vecx)] = \Ex_s[\langle \nabla F(\vecx+s\Delta), \Delta\rangle] / \lambda.
\end{equation*}
\end{restatable}

In Algorihtm~\ref{alg:O2NC-exp-avg}, we set $s_n\sim\mathrm{Exp}(1)$ and then define $\vecx_n = \vecx_{n-1} + s_n\Delta_n$. Thus, Lemma \ref{lem:exp-scaling} implies that 
% $\E[F(\vecx_n) - F(\vecx_{n-1})] = \E\langle \nabla F(\vecx_n), \Delta_n\rangle$. 
\begin{align*}
    \E[F(\vecx_n) - F(\vecx_{n-1})] 
    &= \E\langle \nabla F(\vecx_n), \Delta_n\rangle \\
    &= \E\langle \nabla F(\vecx_n), \vecx_n-\vecx_{n-1}\rangle.
\end{align*}
In other words, we can estimate the ``training progress'' $F(\vecx_n)-F(\vecx_{n-1})$ by directly computing the stochastic gradient at iterate $\vecx_n$. By exploiting favorable properties of the exponential distribution, we dispense with the need for the ``auxiliary point'' $\vecw_n $ employed by O2NC.
% In contrast to the previous scenario where $s_n$ is uniform, this new approach allows us to compute the gradient directly at $\vecx_n$ to estimate the training progress $F(\vecx_n)-F(\vecx_{n-1})$.

We'd like to highlight the significance of this result. The vast majority of smooth non-convex optimization analysis depends on the assumption that $F(\vecx)$ is locally linear, namely $F(\vecx_n)-F(\vecx_{n-1}) \approx \langle \nabla F(\vecx_n), \vecx_n-\vecx_{n-1}\rangle$. Under various smoothness assumptions, the error in this approximation can be controlled via bounds on the remainder in a Taylor series. For example, when $F$ is smooth, then $F(\vecx_n) - F(\vecx_{n-1}) = \langle F(\vecx_n),\vecx_n-\vecx_{n-1}\rangle + O(\|\vecx_n-\vecx_{n-1}\|^2)$. However, since smoothness is necessary for bounding Taylor approximation error, such analysis technique is inapplicable in non-smooth optimization. In contrast, by scaling an exponential random variable to the update, we directly establish a linear equation that $\E[F(\vecx_n)-F(\vecx_{n-1})]=\E\langle \nabla F(\vecx_n),\vecx_n-\vecx_{n-1}\rangle$, effectively eliminating any additional error that Taylor approximation might incur. 

A randomized approach such as ours is also recommended in the recent findings by \citet{jordan2023deterministic}, who demonstrated that randomization is \emph{necessary} for achieving a dimension-free rate in non-smooth optimization. In particular, any deterministic algorithm suffers an additional dimension dependence of $\Omega(d)$. 

Although employing exponential random scaling might seem unconventional, we justify this scaling by noting that $s_n\sim\Expo(1)$ satisfies $\E[s_n]=1$ and $\P\{s_n \ge t\} = \exp(-t)$ (in particular, $\P\{s_n \le 5\} \ge 0.99$). In other words, with high probability, the scaling factor behaves like a constant scaling on the update. To corroborate the efficacy of random scaling, we have conducted a series of empirical tests, the details of which are discussed in Section \ref{sec:experiment}.


\paragraph{Exponentiated and Regularized Losses}
The most significant feature of Exponentiated O2NC (and from which it derives its name) is the loss function: $\ell_n(\Delta) = \langle \beta^{-n}\vecg_n, \Delta\rangle + \calR_n(\Delta)$. This loss consists of two parts: intuitively, the exponentially upweighted linear loss $\langle \beta^{-n}\vecg_n,\Delta\rangle$ measures the ``training progress'' $F(\vecx_n)-F(\vecx_{n-1})$ (as discussed in Lemma \ref{lem:exp-scaling}), and $\calR_n(\Delta)$ serves as an stabilizer that prevents the iterates from irregular behaviors. We will elaborate the role of each component later.
To illustrate how Exponential O2NC works, let $\vecu_n$ be the optimal choice of $\Delta_n$ in hindsight. Then by minimizing the regret $\regret_n(\vecu_n)$ w.r.t. $\vecu_n$, Algorithm \ref{alg:O2NC-exp-avg} automatically chooses the best possible update $\Delta_n$ that is closest to $\vecu_n$.


\paragraph{Exponentially Weighted Gradients}
% The key modification distinguishing Exponentiated O2NC from standard O2NC (and the feature from which it derives its name) is the selection of the loss function. While O2NC uses linear loss $\langle\vecg_n,\Delta\rangle$, we adopt an exponentiated and regularized loss: $\ell_n(\Delta) = \langle \beta^{-n}\vecg_n, \Delta\rangle + \calR_n(\Delta)$. 
For now, set aside the regularizer $\calR_n$ and focus on the linear term $\langle\beta^{-n}\vecg_n,\Delta\rangle$. To provide an intuition why we upweight the gradient by an exponential factor $\beta^{-n}$, we provide a brief overview for the convergence analysis of our algorithm. For simplicity of illustration, we assume $\vecg_n = \nabla F(\vecx_n)$ and $\calR_n=0$.

Let $S_n=\{\vecx_t\}_{t=1}^n$ and let $\vecy_n$ be distributed over $S_n$ such that $\P\{\vecy_n=\vecx_t\} = p_{n,t} := \beta^{n-t}\cdot \frac{1-\beta}{1-\beta^n}$. Our strategy will be to show that this set $S_n$ and random variable $\vecy_n$ satisfy the conditions to make $\overline \vecx_n$ a $(c,\epsilon)$ stationary point  where $\overline\vecx_n$ is defined in Algorithm \ref{alg:O2NC-exp-avg}. To start, note that this distribution satisfies $\overline\vecx_n=\E[\vecy_n]$. Next, since there is always non-zero probability that $\vecy_n=\vecx_1$, it's not possible to obtain a deterministic bound of $\|\vecy_n-\overline\vecx_n\|\le \delta$ for some small $\delta$ (as would be required if we were trying to establish $(\delta,\epsilon)$ Goldstein stationarity). However, since $\vecy_n$ is exponentially more likely to be a later iterate (close to $\vecx_n$) than an early iterate (far from $\vecx_n$), intuitively $\E\|\vecy_n-\overline\vecx_n\|^2$ should not be too big. Formalizing this intuition forms a substantial part of our analysis.

% In fact, we will see in a moment that $\E\|\vecy_n-\overline\vecx_n\|\ll \max\|\vecy_n-\overline\vecx_n\|$. 

In the convergence analysis, we will show $\overline\vecx$ is a $(c,\epsilon)$-stationary point by bounding $\|\nabla F(\overline\vecx_n)\|_c$ (defined in Definition \ref{def:stationary-point}) for all $n$. The condition $\E[\vecy_n]=\overline\vecx_n$ is already satisfied by construction of $\vecy_n$, and it remains to bound the expected gradient $\|\E[\nabla F(\vecy_n)]\|$ and the variance $\E\|\vecy_n-\overline\vecx_n\|^2$. While the regularizer $\calR_n$ is imposed to control the variance, the exponentiated gradients is employed to bound the expected gradient. 
In particular, this is achieved by reducing the difficult task of minimizing the expected gradient of a non-smooth non-convex objective to a relatively easier (and very heavily studied) one: minimizing the regret w.r.t. exponentiated losses $\ell_t(\Delta) = \langle \beta^{-t}\vecg_t,\Delta\rangle$.
To elaborate further, let's consider a simplified illustration as follows.
% To elaborate further, let's first briefly review how O2NC bounds the expected gradient.

% In O2NC (refer to Algorithm \ref{alg:o2nc-original} for notations), \citet{cutkosky2023optimal} draw a connection between the expected gradient and the shifting regret by carefully constructing the comparator. Specifically, by setting $\vecu_n=\vecu^k$ for all $n\in I_k=[(k-1)T+1,kT]$ where
% \begin{equation}
%     \vecu^k = -D \frac{\sum_{t\in I_k} \nabla F(\vecw_n)}{\|\sum_{t\in I_k} \nabla F(\vecw_n)\|},
%     \label{eq:o2nc-u}
% \end{equation}
% it follows that $\vecy'_k\sim \Unif(\{\vecw_n\}_{n\in I_k})$ satisfies
% \begin{align*}
%     \frac{1}{DN} \sum_{n=1}^N \langle \nabla F(\vecw_n), -\vecu_n\rangle 
%     &= \frac{1}{K}\sum_{k=1}^K \left\| \frac{1}{T}\sum_{n\in I_k} \nabla F(\vecw_n) \right\| \\
%     &= \Ex_{k \sim \Unif([K])} \|\E \nabla F(\vecy'_k)\|.
% \end{align*}
% For simplicity of illustration, assume $\vecg_n = \nabla F(\vecw_n)$ and $\sum_{n=1}^N \langle \nabla F(\vecw_n),-\Delta_n\rangle = O(1)$. 
% % (the latter is true with proper assumption because $\E \langle\nabla F(\vecw_n),\Delta_n\rangle = \E[\nabla F(\vecx_n)-F(\vecx_{n-1})]$). 
% Consequently,
% \begin{align*}
%     &\Ex_k\|\E\nabla F(\vecy_k')\| \\
%     &= \frac{1}{DN}\sum_{n=1}^N \langle\nabla F(\vecw_n), \Delta_n - \vecu_n\rangle - \langle \nabla F(\vecw_n), \Delta_n\rangle \\
%     &\lesssim \frac{1}{DN}\regret(\vecu_1,\ldots,\vecu_N).
% \end{align*}
% Here $\regret(\vecu_1,\ldots,\vecu_N)=\sum_{n=1}^N\langle\vecg_n,\Delta_n-\vecu_n\rangle$ denotes the shifting regret with linear loss $\langle \vecg_t,\Delta\rangle$ and comparators $\vecu_n$ defined in \eqref{eq:o2nc-u}. In other words, O2NC reduces the difficult task of bounding the expected gradient to minimizing the shifting regret.

% In our algorithm, such reduction from expected gradient to regret is significantly more complicate since $\vecy_n$ is no longer uniformly distributed. 

% Let $p_{n,t} = P_n(\vecx_t) = \beta^{n-t}\cdot \frac{1-\beta}{1-\beta^n}$. 
Recall that $p_{n,t} = \beta^{n-t}\cdot \frac{1-\beta}{1-\beta^n}$. By construction of $\vecy_n$,
\begin{align*}
    \E[\nabla F(\vecy_n)]
    &= \littlesum_{t=1}^n p_{n,t} \nabla F(\vecx_t).
\end{align*}
Next, for each $n\in[N]$, we define
\begin{align}
    \vecu_n = -D\frac{\sum_{t=1}^n p_{n,t} \nabla F(\vecx_t)}{\|\sum_{t=1}^n p_{n,t}\nabla F(\vecx_t)\|}
    \label{eq:eo2nc-u}
\end{align}
for some $D$ to be specified later. As a remark, $\vecu_n$ minimizes $\langle \E[\nabla F(\vecy_n)], \Delta\rangle$ over all possible $\Delta$ such that $\|\Delta\|=D$, therefore representing the optimal update in iterate $n$ that leads to the fastest convergence.

With $\vecu_n$ defined in \eqref{eq:eo2nc-u}, it follows that
\begin{align*}
    \frac{1}{D}\sum_{t=1}^n p_{n,t} \langle \nabla F(\vecx_t), -\vecu_n\rangle 
    &= \left\| \sum_{t=1}^n p_{n,t} \nabla F(\vecx_t) \right\| \\
    &= \|\E [\nabla F(\vecy_n)]\|.
\end{align*}
Recall that we assume $\vecg_t=\nabla F(\vecx_t)$ for simplicity. Moreover, in later convergence analysis, we will carefully show that $\sum_{n=1}^N\sum_{t=1}^n p_{n,t} \langle\nabla F(\vecx_t),-\Delta_t\rangle \lesssim 1-\beta$ (see Appendix \ref{app:alg}). Consequently,
\begin{align*}
    &\frac{1}{N} \sum_{n=1}^N \|\E\nabla F(\vecy_n)\| \\
    &= \frac{1}{DN} \sum_{n=1}^N \sum_{t=1}^n p_{n,t} \langle \nabla F(\vecx_t), \Delta_t-\vecu_n\rangle \\
    &\quad - \frac{1}{DN} \sum_{n=1}^N \sum_{t=1}^n p_{n,t}\langle \nabla F(\vecx_t),\Delta_t\rangle \\
    &\lesssim \frac{1-\beta}{DN} \left(1+\sum_{n=1}^N \beta^n 
    \regret_n(\vecu_n)\right).
    % \sum_{t=1}^n \langle \beta^{-t}\vecg_t, \Delta_t-\vecu_n\rangle.
\end{align*}
Here $\regret_n(\vecu_n) = \sum_{t=1}^n \langle \beta^{-t}\vecg_t,\Delta_t-\vecu_n\rangle$ denotes the regret w.r.t. the exponentiated losses $\ell_t(\Delta) = \langle \beta^{-t}\vecg_t,\Delta\rangle$ for $t=1,\ldots,n$ (assuming $\calR_n=0$) and comparator $\vecu_n$ defined in \eqref{eq:eo2nc-u}. Notably, the expected gradient is now bounded by the weighted average of the sequence of static regrets, $\regret_n(\vecu_n)$. Consequently, a good OCO algorithm that effectively minimizes the regret is closely aligned with our goal of minimizing the expected gradient.


\paragraph{Variance Regularization}
As aforementioned, we impose the regularizer $\calR_n(\Delta) = \frac{\mu_n}{2}\|\Delta\|^2$ to control the variance $\E\|\vecy_n-\overline\vecx_n\|^2$. Formally, the following result establishes a reduction from bounding variance to bounding the norm of $\Delta_t$, thus motivating the choice of the regularizer.

\begin{restatable}{lemma}{VarianceToRegularizer}
\label{lem:variance-regularizer}
For any $\beta\in(0,1)$, 
\begin{align*}
    \Ex_s\sum_{n=1}^N \Ex_{\vecy_n} \|\vecy_n-\overline\vecx_n\|^2 \le \sum_{n=1}^N \frac{12}{(1-\beta)^2}\|\Delta_n\|^2.
\end{align*}
\end{restatable}

This suggests that bounding $\|\Delta_n\|^2$ is sufficient to bound the variance of $\vecy_n$. Therefore, we impose the regularizer $\calR_n(\Delta) = \frac{\mu_n}{2}\|\Delta\|^2$, for some constant $\mu_n$ to be determined later, to ensure that $\|\Delta_n\|^2$ remains small, effectively controlling the variance of $\vecy_n$. 

Furthermore, we'd like to highlight that Lemma \ref{lem:variance-regularizer} provides a strictly better bound on the variance of $\vecy_n$ compared to the possible maximum deviation $\max\|\vecy_n-\overline\vecx_n\|$. For illustration, assume $\Delta_t$'s are orthonormal, then $\max\|\vecy_N-\overline\vecx_N\| \approx \|\vecx_1-\vecx_N\| = O(N)$. On the other hand, Lemma \ref{lem:variance-regularizer} implies that for $n\sim\Unif([N])$, $\E_n[\Var(\vecy_n)] = O(\frac{1}{(1-\beta)^2})$. In particular, we will show that $1-\beta= N^{-1/2}$ when the objective is smooth. Consequently, $\E\|\vecy_n-\overline\vecx_n\| = O(\sqrt{N})$, which is strictly tighter than the deterministic bound of $\max\|\vecy_N-\overline\vecx_N\| = O(N)$.

This further motivates why we choose this specific distribution for $\vecy_n$: the algorithm does not need to be conservative all the time and can occasionally make relatively large step, breaking the deterministic constraint that $\|\vecy_n-\overline\vecx_n\|\le\delta$, while still satisfying $\Var(\vecy_n)\le\delta^2$.


% \paragraph{Unconstrained OCO}
% EITHER HERE OR LATER IN SEC. 4



\subsection{Convergence Analysis}
\label{sec:alg-conv}

Now we present the main convergence theorem of Algorithm \ref{alg:O2NC-exp-avg}. This is a very general theorem, and we will prove the convergence bound of a more specific algorithm (Theorem \ref{thm:sgdm}) based on this result. A more formally stated version of this theorem and its proof can be found in Appendix \ref{app:alg}.  

\begin{theorem}
\label{thm:o2nc}
Follow Assumption \ref{as:optimization}. Let $\regret_n(\vecu_n)$ denote the regret w.r.t. $\ell_t(\Delta)=\langle\beta^{-t}\vecg_t,\Delta\rangle + \calR_t(\Delta)$ for $t=1,\ldots,n$ and comparator $\vecu_n$ defined in \eqref{eq:eo2nc-u}. Define $\calR_t(\Delta) = \frac{\mu_t}{2}\|\Delta\|^2$, $\mu_t = \frac{24cD}{\alpha^2}\beta^{-t}$, and $\alpha=1-\beta$, then
\begin{align*}
    &\E \|\nabla F(\overline\vecx)\|_c
    \lesssim \frac{F^*}{DN} + \frac{G+\sigma}{\alpha N} + \sigma\sqrt{\alpha} + \frac{c D^2}{\alpha^2} \\
    & + \E\frac{\beta^{N+1}\regret_N(\vecu_N) + \alpha \sum_{n=1}^N \beta^n \regret_n(\vecu_n)}{DN}.
\end{align*}
\end{theorem}

Here the second line denotes the weighted average of the sequence of static regrets, $\regret_n(\vecw_n)$, w.r.t. the exponentiated and regularized loss $\ell_t(\Delta) = \langle \beta^{-t}\vecg_t,\Delta\rangle$ and comparator $\vecu_n$ defined in \eqref{eq:eo2nc-u}, as we discussed earlier. To see an immediate implication of Theorem \ref{thm:o2nc}, assume the average regret is no larger than the terms in the first line. Then by proper tuning $D=\frac{1}{\sqrt{\alpha} N}$ and $\alpha=\max\{\frac{1}{N^{2/3}}, \frac{c^{2/7}}{N^{4/7}}\}$, we have $\E\|\nabla F(\overline\vecw)\|_c \lesssim \frac{1}{N^{1/3}} + \frac{c^{1/7}}{N^{2/7}}$.

% To see a quick example, consider projected online gradient descent (OGD) \cite{zinkevich2003online}, which updates $\Delta_{t+1} = \Pi_D(\Delta_t - \eta\nabla\ell_t(\Delta_t))$ 